% Generated draft: requires mathematical and source review.
\documentclass[11pt]{article}
\usepackage[margin=25mm]{geometry}
\usepackage{fontspec}
\setmainfont{Latin Modern Roman}
\usepackage{amsmath,amssymb,mathtools,unicode-math}
\setmathfont{Latin Modern Math}
\usepackage{enumitem,needspace,xurl,hyperref}
\setlength{\parindent}{0pt}
\setlength{\parskip}{5pt}
\setlength{\emergencystretch}{3em}
\newcommand{\N}{\mathbb{N}}
\newcommand{\Z}{\mathbb{Z}}
\newcommand{\Q}{\mathbb{Q}}
\newcommand{\R}{\mathbb{R}}
\newcommand{\C}{\mathbb{C}}
\newcommand{\F}{\mathbb{F}}
\newcommand{\A}{\mathbb{A}}
\newcommand{\PP}{\mathbb{P}}
\newcommand{\E}{\mathbb{E}}
\newcommand{\eps}{\varepsilon}
\newcommand{\dd}{\,\mathrm{d}}
\begin{document}
\section*{UnsolvedMath ID: 3700036}
\textbf{Source problem code:} catalogue entry
\subsection{Definitions and mathematical statement}
Let \(n\ge1\), let \(a_1,\ldots,a_n>0\) be real charge strengths, and let \(x_1,\ldots,x_n\in\mathbb R^3\) be their locations. On \(\Omega=\mathbb R^3\setminus\{x_1,\ldots,x_n\}\), with Euclidean norm \(|\cdot|\), define the Newtonian potential and its critical set by
\[
 u(x)=\sum_{k=1}^n\frac{a_k}{|x-x_k|},\qquad
 Z(u)=\{x\in\Omega:\nabla u(x)=0\}
 =\left\{x\in\Omega:\sum_{k=1}^n a_k\frac{x-x_k}{|x-x_k|^3}=0\right\}.
\]
The question concerns configurations for which \(Z(u)\) is finite. Its points are counted once each, without a nondegeneracy or multiplicity assumption. What is the largest possible \(|Z(u)|\) in terms of \(n\)? In particular, the proposed Maxwell bound is the assertion
\[
\forall n\ge1\ \forall(a_1,\ldots,a_n)\in(0,\infty)^n\
\forall(x_1,\ldots,x_n)\in(\mathbb R^3)^n:\quad
 |Z(u)|<\infty\ \Longrightarrow\ |Z(u)|\le(n-1)^2.
\]
One may encode the broader extremal question by the supremum of \(|Z(u)|\) over these finite-critical-set configurations, allowing the supremum to be infinite. The conventional point-charge formulation uses distinct locations; the excerpt does not say this explicitly. Coincident positive charges can be combined at their common location, so including coincident locations does not strengthen the proposed universal bound beyond the distinct-location formulation with the appropriate number of locations.

The registry labels the problem solved and its research summary reports a negative resolution in 2026, asserting a five-charge configuration with at least \(24\) critical points, exceeding \((5-1)^2=16\). No citation, configuration, or verification of this report is supplied, so it is retained as a source claim rather than an established counterexample in this draft. The summary also says a sharp asymptotic constant remains under investigation, but does not specify the normalization or an exact asymptotic conjecture. That broader issue is distinguished from the yes/no Maxwell-bound question.
\subsection{Short English statement}
How many isolated equilibrium points can the Newtonian potential of \(n\) positive charges in \(\mathbb R^3\) have, assuming its entire critical set is finite? The proposed upper bound is \((n-1)^2\). The supplied record reports a 2026 counterexample but provides no supporting citation, so its solved label requires source review.
\subsection{Sources}
\begin{itemize}
\item Local registry: \url{data/unsolvedmath/problems.json}, permanent UnsolvedMath ID \(3700036\), source code catalogue entry; supplied statement, background, and research summary, updated 2026-08-06. Attribution: Alexandre Eremenko and cited proposers; registry proposed year 2008.
\item Alexandre Eremenko, \emph{My favorite unsolved problems}, ``Maxwell problem'', Question 2: \url{https://www.math.purdue.edu/~eremenko/uns1.html}. Supplied extraction method: PDF text; accessed 2026-07-29. The record reports applying explicit solved/counterexample updates in the author-maintained index through that access date.
\item Maxwell's bound and a purported 2026 five-charge counterexample are mentioned in the supplied record, but no original Maxwell reference or identifiable counterexample publication is supplied.
\end{itemize}
\end{document}
